\documentclass[12pt, a4paper]{article}

% ── Packages ──────────────────────────────────────────────────────────────────
\usepackage[utf8]{inputenc}
\usepackage[T1]{fontenc}
\usepackage{lmodern}
\usepackage[margin=1in]{geometry}
\usepackage{graphicx}
\usepackage{amsmath, amssymb}
\usepackage{booktabs}
\usepackage{hyperref}
\usepackage{xcolor}
\usepackage{enumitem}
\usepackage{float}
\usepackage{caption}
\usepackage{parskip}
\usepackage{tikz}
\usetikzlibrary{positioning, arrows.meta, shapes.geometric}
\usepackage{fancyhdr}
\usepackage{multicol}

% ── Colours ───────────────────────────────────────────────────────────────────
\definecolor{vibeblue}{RGB}{70,130,180}
\definecolor{vibeorange}{RGB}{230,126,34}
\definecolor{vibegreen}{RGB}{39,174,96}
\definecolor{vibepurple}{RGB}{142,68,173}
\definecolor{vibered}{RGB}{192,57,43}

% ── Header / Footer ──────────────────────────────────────────────────────────
\pagestyle{fancy}
\fancyhf{}
\rhead{Beyond Play Count --- Advanced Analysis Possibilities}
\lhead{ML Project --- Exploration Report}
\cfoot{\thepage}

% ── Title ─────────────────────────────────────────────────────────────────────
\title{
    \vspace{-1cm}
    \textbf{Beyond Play Count:\\
    Exploring Advanced Time Series Analyses\\
    and ``Listening Vibes'' for Music Data}\\[0.5em]
    \large An In-Depth Exploration of Possibilities, Metrics,\\
    and Machine Learning Approaches Using \texttt{fm\_data.csv}
}
\author{ML Project Team}
\date{\today}

% ══════════════════════════════════════════════════════════════════════════════
\begin{document}
\maketitle
\thispagestyle{fancy}

\begin{abstract}
The initial project report established a baseline approach: aggregating
Last.fm scrobble data into daily play counts and forecasting them with
SARIMA. While effective as a starting point, raw play count is the
\emph{least informative} metric we can extract from the data. This document
explores a rich landscape of alternative analyses --- from defining temporal
``listening vibes'' (clusters of songs based on when they were played) to
computing diversity indices, discovery rates, loyalty scores, and session
dynamics. For each possibility, we discuss what the metric captures, how it
would be constructed from \texttt{fm\_data.csv}, what time series plots it
enables, and which machine learning models are most suitable for forecasting
or classification. No code is included; the focus is purely on conceptual
exploration and modelling strategy.
\end{abstract}

\tableofcontents
\newpage

% ══════════════════════════════════════════════════════════════════════════════
\section{Motivation: Why Move Beyond Play Count?}
% ══════════════════════════════════════════════════════════════════════════════

Daily play count answers a single question: \emph{How much did the user
listen?} But it says nothing about:

\begin{itemize}
    \item \textbf{When} they listened (morning vs.\ late night),
    \item \textbf{What} they listened to (one artist on repeat vs.\ dozens),
    \item \textbf{How} they listened (short background sessions vs.\ long
          focused ones),
    \item \textbf{Whether} they were exploring new music or revisiting
          favourites.
\end{itemize}

Our dataset contains \textbf{Username, Artist, Track, Album, Date, and Time}
for 166{,}153 scrobbles across 11 users over 31 days. This is far richer
than a single count. The following sections explore every meaningful metric
and analysis we can extract from it.

% ══════════════════════════════════════════════════════════════════════════════
\section{The ``Listening Vibes'' Framework}
% ══════════════════════════════════════════════════════════════════════════════

\subsection{What Are Vibes?}

A \textbf{vibe} is a temporal cluster --- a named time window defined by
when a scrobble occurred. The hypothesis is simple: \emph{people listen to
different kinds of music at different times, and these temporal slots carry
distinct emotional and behavioural signatures.}

Rather than treating time as a raw number (hour 0--23), we assign each
scrobble to a meaningful, human-interpretable category.

\subsection{Proposed Vibe Definitions}

We propose two tiers of vibes: \textbf{primary} (based on time of day) and
\textbf{compound} (combining day of week with time of day).

\subsubsection{Primary Vibes (Time-of-Day)}

\begin{center}
\begin{tabular}{llp{6.5cm}}
    \toprule
    \textbf{Vibe} & \textbf{Hours} & \textbf{Behavioural Hypothesis} \\
    \midrule
    Early Morning  & 05:00--08:59 & Wake-up routines, commute music.
                                    Likely energetic or habitual. \\
    Morning        & 09:00--11:59 & Focus and productivity.
                                    Instrumental, ambient, or familiar. \\
    Afternoon      & 12:00--16:59 & Midday energy. Mixed genres,
                                    background listening. \\
    Evening        & 17:00--20:59 & Wind-down or social.
                                    More intentional, curated listening. \\
    Late Night     & 21:00--00:59 & Reflective, moody, or intense.
                                    Deeper engagement. \\
    Deep Night     & 01:00--04:59 & Insomnia or marathon sessions.
                                    Niche, experimental. \\
    \bottomrule
\end{tabular}
\end{center}

\subsubsection{Compound Vibes (Day + Time)}

\begin{center}
\begin{tabular}{llp{6.5cm}}
    \toprule
    \textbf{Vibe} & \textbf{Definition} & \textbf{Behavioural Hypothesis} \\
    \midrule
    Weekend Vibes   & Sat--Sun, any time  & Leisure listening, longer
                                            sessions, more exploration. \\
    Weekday Grind   & Mon--Fri, 09--17    & Work-time music, focus-oriented,
                                            lower diversity. \\
    Friday Night     & Fri 18:00 -- Sat 03:00 & Social, party, high-energy
                                                or celebratory. \\
    Sunday Chill     & Sun, all day        & Recovery, mellow, singer-songwriter,
                                            acoustic. \\
    Monday Morning   & Mon, 05:00--11:59   & Re-entry to the work week,
                                            motivational or comforting. \\
    \bottomrule
\end{tabular}
\end{center}

\subsection{What Makes Vibes Analytically Powerful?}

Vibes are not arbitrary labels --- they are \textbf{categorical features}
that enable a fundamentally new type of analysis. Once every scrobble is
tagged with a vibe, we can:

\begin{enumerate}
    \item \textbf{Compare distributions}: Which artists or tracks dominate
          each vibe? Do Late Night and Morning vibes share any overlap?
    \item \textbf{Track vibe proportions over time}: Plot the percentage of
          each vibe per week as a stacked area chart. This reveals
          behavioural shifts (e.g.\ a user gradually becoming a night
          listener).
    \item \textbf{Cross vibes with other metrics}: Compute diversity,
          discovery rate, or loyalty score \emph{within each vibe
          separately}. This is far more revealing than global metrics.
    \item \textbf{Use vibes as features in ML models}: Pass vibe proportions
          as exogenous variables to forecasting models.
\end{enumerate}

\subsection{Vibes as a Time Series}

\subsubsection{Vibe Composition Plot (Stacked Area)}
For each time unit (day or week), compute the proportion of scrobbles falling
into each vibe. Plotted as a \textbf{stacked area chart}, this shows how the
user's daily rhythm evolves over the 31-day window.

\textbf{What to look for}:
\begin{itemize}
    \item Do weekends consistently show a different vibe mix?
    \item Are there days where Late Night and Deep Night vibes spike
          (suggesting binge sessions)?
    \item Does the vibe composition remain stable, or does it drift over
          the month?
\end{itemize}

\subsubsection{Vibe Volume Plot (Absolute Counts)}
Instead of proportions, plot the raw number of scrobbles per vibe per day.
This combines the ``how much'' of play count with the ``when'' of vibes,
giving a much richer picture.

\subsubsection{Machine Learning Approach for Vibes}

\paragraph{Forecasting vibe proportions.}
Each day produces a vector of vibe proportions
$\mathbf{v}_t = (v_t^{\text{morning}}, v_t^{\text{afternoon}}, \ldots)$.
This is a \textbf{compositional time series} (proportions that sum to 1).

\begin{itemize}
    \item \textbf{Vector Autoregression (VAR)}: Treats each vibe proportion
          as one variable in a multivariate time series. VAR captures
          cross-correlations --- e.g.\ a spike in Late Night today might
          predict a drop in Early Morning tomorrow.
    \item \textbf{Dirichlet Regression}: Since vibes are compositional
          (proportions summing to 1), Dirichlet regression is the
          theoretically correct model. It respects the simplex constraint
          and avoids the problems of modelling proportions independently.
    \item \textbf{LSTM / GRU}: A recurrent neural network that takes the
          vibe vector $\mathbf{v}_{t-k}, \ldots, \mathbf{v}_{t-1}$ as input
          and predicts $\mathbf{v}_t$. Captures non-linear temporal
          dependencies. Suitable if the dataset is extended beyond 31 days.
\end{itemize}

\paragraph{Classifying vibes from artist/track features.}
Given a scrobble's artist and track (but \emph{not} its timestamp), can we
predict which vibe it belongs to? This asks: \emph{Is there a systematic
relationship between what someone listens to and when they listen?}

\begin{itemize}
    \item \textbf{Features}: Artist frequency encoding, track repeat count,
          album dummy variables, or artist embeddings from SVD.
    \item \textbf{Models}: Multinomial logistic regression (baseline),
          random forest, gradient-boosted trees (XGBoost/LightGBM).
    \item \textbf{Evaluation}: Multi-class accuracy, macro F1 score,
          confusion matrix between vibes.
    \item \textbf{Insight}: If accuracy is significantly above random
          ($\sim$16--20\% for 5--6 vibes), this proves that listening
          content is predictive of listening time --- people do have
          ``morning artists'' and ``night artists.''
\end{itemize}


% ══════════════════════════════════════════════════════════════════════════════
\section{Listening Diversity Index (Shannon Entropy)}
% ══════════════════════════════════════════════════════════════════════════════

\subsection{What It Measures}

The \textbf{Shannon entropy} of a user's listening distribution measures how
\emph{varied} or \emph{uniform} their artist choices are within a given
time window. It answers: \emph{Is the user exploring many artists equally,
or obsessively replaying one?}

\[
    H(t) = -\sum_{i=1}^{N_t} p_i(t) \log_2 p_i(t)
\]

where $p_i(t)$ is the proportion of plays for artist $i$ in time window $t$,
and $N_t$ is the number of distinct artists.

\begin{itemize}
    \item $H = 0$: The user listened to exactly one artist (zero diversity).
    \item $H = \log_2 N$: The user listened to all $N$ artists equally
          (maximum diversity).
\end{itemize}

\subsection{Normalized Entropy}

To compare across time windows with different numbers of artists, we
normalise:
\[
    H_{\text{norm}}(t) = \frac{H(t)}{\log_2 N_t} \in [0, 1]
\]

A value of 0.8 means the user is using 80\% of the available diversity.

\subsection{Time Series Plot}

Compute $H(t)$ weekly (or daily) and plot it over the 31-day window. This
reveals:
\begin{itemize}
    \item \textbf{Exploration phases}: Periods where $H$ rises sharply ---
          the user is sampling many new artists.
    \item \textbf{Obsession phases}: Periods where $H$ drops --- the user
          found something they love and is playing it on repeat.
    \item \textbf{Weekday/weekend patterns}: Diversity might systematically
          differ between work days and free days.
\end{itemize}

\subsection{Cross with Vibes}

Compute entropy \emph{separately for each vibe}. This answers questions like:
\begin{itemize}
    \item Are Late Night sessions more diverse than Morning sessions?
    \item Do Weekend Vibes have higher entropy than Weekday Grind?
    \item Does the user always listen to the same artist at a particular
          time, but explore at another?
\end{itemize}

\subsection{Machine Learning Approach}

\paragraph{Forecasting diversity.}
$H(t)$ is a univariate continuous time series --- any standard time series
model applies.

\begin{itemize}
    \item \textbf{SARIMA}: Directly applicable. Seasonal period of 7 if
          weekly cycles exist in diversity.
    \item \textbf{SARIMAX with exogenous variables}: Use play count,
          vibe proportions, and day-of-week as exogenous regressors to
          predict diversity. This asks: \emph{Can we forecast how varied
          someone's listening will be based on when and how much they
          listen?}
    \item \textbf{Random Forest Regressor}: For non-linear relationships.
          Input features: lagged entropy values, day of week, total plays,
          number of unique artists in the previous $k$ days.
\end{itemize}

\paragraph{Anomaly detection.}
Sudden drops in diversity (obsession events) or spikes (exploration bursts)
are anomalies. An \textbf{Isolation Forest} or \textbf{autoencoder} trained
on normal diversity patterns could flag unusual listening behaviour.

% ══════════════════════════════════════════════════════════════════════════════
\section{New Artist Discovery Rate}
% ══════════════════════════════════════════════════════════════════════════════

\subsection{What It Measures}

For each scrobble, we can check whether the artist has \emph{ever been
played before} by that user (up to that point in time). This yields two
metrics:

\begin{enumerate}
    \item \textbf{New artists per day/week}: The absolute count of
          first-ever artist listens within the time window.
    \item \textbf{Discovery rate}: The fraction of scrobbles in a time
          window that feature a never-before-heard artist.
\end{enumerate}

\subsection{Cumulative Discovery Curve}

The running total of unique artists over time follows a characteristic shape:
\begin{itemize}
    \item \textbf{Early}: Steep rise (most artists are ``new'').
    \item \textbf{Later}: Flattening curve (the user has heard most of their
          rotation).
\end{itemize}

The \emph{rate of flattening} reveals whether the user's taste is saturating
or still actively growing. Over a longer dataset, this curve typically
follows \textbf{Heaps' law}: $V(n) \approx K \cdot n^\beta$ where
$V(n)$ is the vocabulary (unique artists) after $n$ plays, and $\beta < 1$.

\subsection{Time Series Plot}

Plot the daily discovery rate as a line chart with a rolling average overlay.
Spikes indicate days where the user actively sought out new music; troughs
indicate comfort-listening days.

\subsection{Cross with Vibes}

Compute the discovery rate separately per vibe. This answers:
\begin{itemize}
    \item \emph{Do users discover new music more during Late Night sessions
          or during Weekend Vibes?}
    \item \emph{Is the Weekday Grind purely repetitive, or does it also
          introduce new artists?}
\end{itemize}

\subsection{Machine Learning Approach}

\paragraph{Forecasting discovery rate.}
This is another univariate time series, but with a particular characteristic:
it tends to \textbf{decrease over time} as the user accumulates more known
artists.

\begin{itemize}
    \item \textbf{SARIMA with trend}: The $d$ parameter in ARIMA handles
          the decreasing trend. A model of $(1,1,1)(1,0,1,7)$ could work.
    \item \textbf{Exponential smoothing (ETS)}: Suitable for data with a
          decaying trend. A damped trend model (ETS with
          $\text{trend}=\text{damped}$) would naturally capture the
          flattening behaviour.
    \item \textbf{Poisson regression}: Since ``new artists per day'' is a
          count, a Poisson or negative binomial regression with time and
          vibe features as regressors is statistically appropriate.
\end{itemize}

\paragraph{Classification: Will the user discover a new artist today?}
Binary: yes/no, at least one new artist. Features: day of week, vibe
distribution of previous days, recent entropy, cumulative unique count so
far. Model: logistic regression or gradient-boosted classifier.

% ══════════════════════════════════════════════════════════════════════════════
\section{Artist Rank Dynamics (Personal Billboard Chart)}
% ══════════════════════════════════════════════════════════════════════════════

\subsection{What It Measures}

For each time window (day or week), rank the user's top $N$ artists by play
count. Track how these ranks change over time --- like a personal Billboard
Hot 100.

\subsection{Time Series Plot}

A \textbf{bump chart} (ranked line chart with inverted $y$-axis: rank 1 at
top). Each artist is a coloured line. Lines crossing each other show
taste evolution.

\textbf{What to look for}:
\begin{itemize}
    \item \textbf{Stable monarchs}: Artists that hold rank 1 consistently.
    \item \textbf{Meteoric rises}: An artist suddenly jumping from unranked
          to top 3 (a new obsession).
    \item \textbf{Fadeouts}: Artists dropping out of the top $N$ over time
          (novelty wearing off).
\end{itemize}

\subsection{Machine Learning Approach}

\paragraph{Rank forecasting.}
Predicting an artist's rank next week given their recent trajectory.

\begin{itemize}
    \item \textbf{Ordinal regression}: Ranks are ordinal --- standard
          regression treats the difference between rank 1 and 2 the same as
          between rank 10 and 11, which is not ideal. Ordinal regression
          preserves the ordering without assuming equal spacing.
    \item \textbf{Learning-to-rank (LTR)}: Algorithms like
          \textbf{LambdaMART} (used in search engine ranking) can learn to
          predict artist rankings from features like recent play counts,
          trajectory direction, and consistency.
    \item \textbf{Per-artist ARIMA}: Fit a separate ARIMA model to each
          artist's daily play count, then rank the resulting forecasts. This
          is a simple but effective decomposition.
\end{itemize}

\paragraph{Detecting rank change events.}
A \textbf{change-point detection} algorithm (e.g.\ PELT, BOCPD) applied to
each artist's play count series can identify the exact day a new obsession
began or an old favourite faded.

% ══════════════════════════════════════════════════════════════════════════════
\section{Listening Session Analysis}
% ══════════════════════════════════════════════════════════════════════════════

\subsection{Defining a Session}

A \textbf{listening session} is a continuous period of listening activity.
We define a session boundary as a gap of more than 30 minutes between
consecutive scrobbles. Every scrobble within a session belongs to the same
session ID.

\subsection{Session-Level Metrics}

From each session we can extract:

\begin{center}
\begin{tabular}{lp{8.5cm}}
    \toprule
    \textbf{Metric} & \textbf{Description} \\
    \midrule
    Duration (minutes) & Time from first to last scrobble in the session. \\
    Track count        & Number of songs played in the session. \\
    Unique artists     & Number of distinct artists in the session. \\
    Artist entropy     & Shannon diversity within the session. \\
    Repeat ratio       & Fraction of tracks that are repeats within
                         the session. \\
    \bottomrule
\end{tabular}
\end{center}

\subsection{Time Series Plots}

Aggregate session-level metrics daily or weekly:
\begin{itemize}
    \item \textbf{Average session duration over time}: Reveals whether the
          user is engaging in longer or shorter listening blocks.
    \item \textbf{Sessions per day}: Distinguishes between ``many short
          sessions'' and ``few long sessions'' --- both could produce the
          same raw play count but indicate very different behaviour.
    \item \textbf{Average tracks per session}: A decreasing trend might
          indicate declining engagement; an increasing trend suggests
          deeper listening.
    \item \textbf{Average unique artists per session}: Measures within-session
          exploration. A session with 20 tracks from 1 artist is very
          different from 20 tracks from 15 artists.
\end{itemize}

\subsection{Machine Learning Approach}

\paragraph{Forecasting session characteristics.}
Each daily aggregate (e.g.\ mean session duration) is a time series that can
be modelled independently.

\begin{itemize}
    \item \textbf{SARIMA / ETS}: For the aggregate metrics (mean duration,
          session count per day).
    \item \textbf{Multivariate VAR}: Model all session metrics jointly.
          Session count, mean duration, and mean diversity are likely
          correlated --- a day with many sessions might have shorter average
          duration.
    \item \textbf{Gaussian Process Regression}: For smooth, uncertainty-aware
          forecasts of session metrics with limited data (31 days). GPs
          naturally handle small datasets and provide confidence intervals.
\end{itemize}

\paragraph{Session type classification.}
Cluster sessions into types: ``background listening'' (long, low diversity,
one artist), ``active exploration'' (many unique artists, high entropy),
``party mode'' (Friday night, high energy). Then track how the proportion
of session types changes over time --- this is itself a compositional
time series (similar to vibe proportions).

\begin{itemize}
    \item \textbf{K-Means / Gaussian Mixture Models}: On the feature vector
          (duration, track count, unique artists, entropy, repeat ratio) to
          discover natural session types.
    \item \textbf{Hidden Markov Model (HMM)}: Model the user as transitioning
          between latent ``listening states'' (background, active, binge).
          The HMM captures the temporal dynamics of transitions between
          session types --- e.g.\ a binge session is more likely to follow
          another binge session than to follow a background session.
\end{itemize}

% ══════════════════════════════════════════════════════════════════════════════
\section{Loyalty Score}
% ══════════════════════════════════════════════════════════════════════════════

\subsection{What It Measures}

The \textbf{loyalty score} is the percentage of a user's plays in a time
window that go to their all-time top $K$ artists (e.g.\ $K = 10$):

\[
    L(t) = \frac{\text{plays of top-}K\text{ artists in window } t}
                {\text{total plays in window } t} \times 100
\]

\begin{itemize}
    \item $L = 100\%$: The user exclusively listens to their favourites.
    \item $L \approx 0\%$: The user is entirely outside their usual rotation.
\end{itemize}

\subsection{Time Series Plot}

Plot $L(t)$ daily with a filled area chart. Overlay a horizontal reference
line at 50\%. Periods above 50\% are ``loyal'' phases; periods below are
``adventurous'' phases.

\subsection{Relationship to Diversity}

Loyalty and diversity are \textbf{inversely related} but not identical.
A user could have low loyalty (rarely playing top-10 artists) but also low
diversity (obsessing over one \emph{new} artist). Together, they paint a
more complete picture.

\subsection{Machine Learning Approach}

\begin{itemize}
    \item \textbf{SARIMA}: Loyalty is bounded (0--100\%) but often stays
          in the 20--80\% range, so SARIMA works well in practice.
    \item \textbf{Beta regression}: Since $L(t)/100 \in (0,1)$ is a
          proportion, beta regression is the statistically appropriate
          model. Predictors: day of week, recent diversity, recent session
          count.
    \item \textbf{Threshold classification}: Define phases as ``loyal''
          ($L > 60\%$) or ``exploratory'' ($L \leq 60\%$) and train a
          classifier to predict tomorrow's phase from recent features.
          Model: logistic regression, SVM, or a simple decision tree.
\end{itemize}

% ══════════════════════════════════════════════════════════════════════════════
\section{Seasonal and Cyclical Patterns}
% ══════════════════════════════════════════════════════════════════════════════

\subsection{Weekly Cycles}

The most prominent cycle in listening data is the \textbf{weekly cycle}.
With our data spanning exactly 31 days (approximately 4.4 weeks), we can
detect weekly patterns by:

\begin{enumerate}
    \item \textbf{Day-of-week aggregation}: Average play count (and all
          other metrics) by day of week across all weeks. Visualise as a
          bar chart.
    \item \textbf{Fourier analysis}: Compute the periodogram or FFT of the
          daily time series to identify dominant frequencies. A peak at
          frequency $f = 1/7$ confirms weekly seasonality.
\end{enumerate}

\subsection{Artist Seasonality Heatmap}

For each of the top $N$ artists, compute total plays per day of week
(aggregated across all weeks). Display as a heatmap with artists on the
$y$-axis and days of week on the $x$-axis.

\textbf{What to look for}:
\begin{itemize}
    \item Artists that appear exclusively on weekends.
    \item Artists that spike on specific days (e.g.\ ``Monday motivation''
          artists).
    \item Uniform artists that are listened to equally every day (true
          favourites vs.\ context-dependent artists).
\end{itemize}

\subsection{Hourly Distribution (Listening Clock)}

A \textbf{radial (polar) bar chart} showing play count by hour of day,
plotted around a 24-hour clock. Optionally split into two overlaid rings:
weekday vs.\ weekend. This is not a time series per se, but it
\emph{aggregates} the time series into a powerful visual summary.

\subsection{Machine Learning Approach}

\paragraph{Fourier features for SARIMAX.}
Instead of relying solely on SARIMA's seasonal component, we can explicitly
construct Fourier features:
\[
    \sin\!\left(\frac{2\pi k t}{7}\right), \quad
    \cos\!\left(\frac{2\pi k t}{7}\right)
    \qquad k = 1, 2, 3
\]

These sine/cosine pairs, passed as exogenous variables to SARIMAX, capture
weekly patterns more flexibly than a single seasonal period.

\paragraph{Facebook Prophet.}
Prophet is designed specifically for time series with strong seasonal effects
and multiple seasonality levels. It would model:
\begin{itemize}
    \item \textbf{Weekly seasonality}: Automatic weekday/weekend patterns.
    \item \textbf{Trend}: Piecewise linear trend over the month.
    \item \textbf{Holidays/events}: Could mark specific dates (e.g.\ New
          Year's Day falls within our January 2021 data) as special events.
\end{itemize}

% ══════════════════════════════════════════════════════════════════════════════
\section{Track Repeat Intensity}
% ══════════════════════════════════════════════════════════════════════════════

\subsection{What It Measures}

How many times the user replays the \emph{same track} within a day or
session. This captures the phenomenon of ``having a song on repeat'' ---
a strong signal of emotional connection.

\[
    R(t) = \max_{\text{track } j}\left(\text{play count of } j
           \text{ on day } t\right)
\]

Alternatively, the \textbf{average repeat count}: mean number of times each
played track is repeated within the day.

\subsection{Time Series Plot}

Plot $R(t)$ daily. Spikes indicate days with obsessive replays. Cross with
the identity of the track being repeated to tell a ``story'' of obsessions
over time.

\subsection{Machine Learning Approach}

\begin{itemize}
    \item \textbf{Poisson / Negative Binomial regression}: Repeat counts
          are non-negative integers with possible overdispersion. A count
          model with day-of-week and vibe features as regressors is natural.
    \item \textbf{Zero-inflated models}: Many tracks are played exactly once.
          A zero-inflated Poisson model handles the excess of ``no repeats''
          entries.
    \item \textbf{Survival analysis}: Model the ``time until an artist is
          replayed'' using Cox proportional hazards or Kaplan--Meier
          estimates. Artists with short replay times are ``sticky'';
          artists with long replay times were one-hit curiosities.
\end{itemize}

% ══════════════════════════════════════════════════════════════════════════════
\section{Artist Lifespan Analysis}
% ══════════════════════════════════════════════════════════════════════════════

\subsection{What It Measures}

For each artist, compute:
\begin{itemize}
    \item \textbf{First listen date}: The earliest scrobble of that artist.
    \item \textbf{Last listen date}: The most recent scrobble.
    \item \textbf{Active span}: Last $-$ First (in days).
    \item \textbf{Total plays within span}: How intensely the artist was
          played during their active period.
\end{itemize}

\subsection{Visualisation}

A \textbf{Gantt-style chart} with artists on the $y$-axis and time on the
$x$-axis. Each artist gets a horizontal bar from first to last listen, with
colour or thickness encoding play intensity. This shows which artists
``stuck'' and which were fleeting encounters.

\subsection{Machine Learning Approach}

\paragraph{Predicting artist longevity.}
Given the features of an artist's first few listens (time of day, vibe,
number of initial plays, whether it was discovered via a binge session),
can we predict whether the artist will still be in rotation a week later?

\begin{itemize}
    \item \textbf{Survival models}: Cox regression with features from the
          first session. The ``event'' is the artist being dropped (no plays
          for $> k$ days).
    \item \textbf{Binary classification}: Will this artist be played again
          after 7 days? Features: initial play intensity, vibe of discovery,
          session type, diversity at time of discovery.
    \item \textbf{Random forest / XGBoost}: For the binary classifier, with
          feature importance analysis revealing what predicts artist
          stickiness.
\end{itemize}

% ══════════════════════════════════════════════════════════════════════════════
\section{Listening Gap Analysis}
% ══════════════════════════════════════════════════════════════════════════════

\subsection{What It Measures}

Periods with \textbf{zero or near-zero scrobbles}: the user stopped
listening. These gaps correspond to real-life events --- sleep, work, travel,
social activities, or simply not being in the mood for music.

\subsection{Time Series Plot}

Overlay ``gap regions'' (shaded bands) on top of the standard daily play
count chart. Alternatively, plot the \textbf{daily maximum gap duration}
(longest silence in hours) as its own time series.

\subsection{Machine Learning Approach}

\begin{itemize}
    \item \textbf{Anomaly detection}: Unusually long gaps within a period
          where the user is normally active could flag interesting behavioural
          changes. Isolation Forest or DBSCAN on gap durations.
    \item \textbf{Predicting active periods}: Given the time of day and day
          of week, predict whether the user will be listening. This is a
          binary time series classification problem solvable with logistic
          regression or an HMM with ``active'' and ``inactive'' states.
\end{itemize}

% ══════════════════════════════════════════════════════════════════════════════
\section{Rolling Unique Artist Count}
% ══════════════════════════════════════════════════════════════════════════════

\subsection{What It Measures}

A \textbf{rolling window} (e.g.\ 7 days) count of unique artists. Unlike
instantaneous diversity (entropy), this captures the breadth of the user's
repertoire over a sustained period. It smooths out daily fluctuations and
reveals the underlying trend.

\subsection{Time Series Plot}

A smooth line chart of the 7-day rolling unique artist count. Compare it
against the rolling play count to see if the user is listening more
\emph{broadly} or just listening \emph{more}.

\subsection{Machine Learning Approach}

\begin{itemize}
    \item \textbf{SARIMA}: Applicable to the smoothed series.
    \item \textbf{Regression}: Predict the rolling unique count from
          upcoming features (is a holiday approaching? Is it a weekend
          streak?).
\end{itemize}

% ══════════════════════════════════════════════════════════════════════════════
\section{Multi-User Comparative Analysis}
% ══════════════════════════════════════════════════════════════════════════════

With 11 users in the dataset, every metric above can be computed per user
and compared.

\subsection{Possible Comparisons}

\begin{center}
\begin{tabular}{lp{8cm}}
    \toprule
    \textbf{Comparison} & \textbf{Insight} \\
    \midrule
    Diversity trajectories & Which users explore more over time? Who is
                             most set in their ways? \\
    Vibe profiles         & Do some users listen exclusively at night while
                             others are morning listeners? \\
    Discovery curves      & Whose taste is saturating fastest? Who is the
                             most adventurous? \\
    Loyalty rankings      & Who is the most loyal listener? Who has the
                             most scattered taste? \\
    Session patterns      & Who listens in long focused sessions vs.\ many
                             short bursts? \\
    \bottomrule
\end{tabular}
\end{center}

\subsection{Machine Learning Approach}

\paragraph{User clustering.}
Compute a feature vector for each user: average diversity, average loyalty,
dominant vibe, discovery rate, session characteristics. Apply
\textbf{K-Means} or \textbf{hierarchical clustering} to group users into
``listener types.''

\paragraph{User-level forecasting.}
Fit a \textbf{hierarchical (multi-level) time series model} across all
users simultaneously. Global patterns (everyone listens less on Mondays)
inform individual forecasts, and individual-level parameters capture
user-specific behaviour. This approach can be implemented with
\textbf{mixed-effects models} or \textbf{global neural networks} (one LSTM
trained on all users with user embeddings).

% ══════════════════════════════════════════════════════════════════════════════
\section{Comprehensive Summary of All Possibilities}
% ══════════════════════════════════════════════════════════════════════════════

The following table summarises every metric, its time series potential,
and the recommended machine learning approach:

\begin{center}
\small
\begin{tabular}{p{3.2cm}p{3cm}p{3.5cm}p{4cm}}
    \toprule
    \textbf{Metric} & \textbf{Time Series Type} & \textbf{Best Plot} &
    \textbf{Recommended ML Model} \\
    \midrule
    Daily Play Count
        & Univariate, integer
        & Line chart
        & SARIMA (baseline) \\[4pt]
    Vibe Proportions
        & Multivariate, compositional
        & Stacked area chart
        & VAR, Dirichlet regression, LSTM \\[4pt]
    Shannon Entropy (Diversity)
        & Univariate, continuous
        & Line with filled area
        & SARIMAX, Random Forest, GP \\[4pt]
    Discovery Rate
        & Univariate, decaying
        & Line + rolling avg
        & ETS (damped trend), Poisson regression \\[4pt]
    Artist Rankings
        & Multivariate, ordinal
        & Bump chart
        & Ordinal regression, LambdaMART \\[4pt]
    Session Metrics
        & Multivariate, continuous
        & 2$\times$2 subplot panel
        & VAR, Gaussian Process \\[4pt]
    Session Types
        & Compositional
        & Stacked bar chart
        & K-Means + HMM \\[4pt]
    Loyalty Score
        & Univariate, bounded
        & Filled line chart
        & Beta regression, SARIMA \\[4pt]
    Seasonal Patterns
        & Periodic
        & Heatmap, polar chart
        & Prophet, Fourier + SARIMAX \\[4pt]
    Repeat Intensity
        & Univariate, count
        & Bar chart with labels
        & Poisson / NB regression \\[4pt]
    Artist Lifespan
        & Non-temporal (cross-sectional)
        & Gantt chart
        & Survival analysis (Cox) \\[4pt]
    Listening Gaps
        & Univariate, duration
        & Shaded overlay
        & Isolation Forest, HMM \\[4pt]
    Rolling Unique Artists
        & Univariate, smooth
        & Smooth line chart
        & SARIMA, regression \\[4pt]
    User Clusters
        & Cross-sectional
        & Scatter / radar chart
        & K-Means, hierarchical \\
    \bottomrule
\end{tabular}
\end{center}

% ══════════════════════════════════════════════════════════════════════════════
\section{Recommended Priority Order}
% ══════════════════════════════════════════════════════════════════════════════

Given the constraints of the dataset (31 days, 11 users), some analyses are
more feasible and impactful than others. We recommend the following priority:

\begin{enumerate}[label=\textbf{\arabic*.}]
    \item \textbf{Vibe Composition Over Time} --- The most novel and visually
          compelling analysis. Directly builds on existing feature
          engineering. Enables all downstream cross-analyses.

    \item \textbf{Shannon Diversity Index} --- Conceptually rich, easy to
          compute, and produces a beautiful time series that tells a
          psychological story about each user.

    \item \textbf{Discovery Rate} --- Natural complement to diversity.
          Together, diversity + discovery fully characterise the
          ``exploration vs.\ exploitation'' behaviour.

    \item \textbf{Session Analysis} --- Adds a dimension no other metric
          captures: \emph{how} the user engages (deep sessions vs.\
          scattered plays).

    \item \textbf{Loyalty Score} --- Quick to compute, intuitive, and
          works as a high-level summary of taste stability.

    \item \textbf{Artist Rank Chart} --- Visually distinctive (bump chart),
          tells individual artist ``stories,'' and is immediately
          understandable.

    \item \textbf{Seasonal Heatmap and Listening Clock} --- Best for
          presentation slides. Not a time series per se, but essential for
          understanding the data.

    \item \textbf{Track Repeat Intensity and Artist Lifespan} --- More
          specialised but still interesting. Implement if time allows.

    \item \textbf{Multi-User Comparison} --- A natural final chapter,
          applying all the above across the 11 users to find listener
          archetypes.
\end{enumerate}

% ══════════════════════════════════════════════════════════════════════════════
\section{Constraints and Considerations}
% ══════════════════════════════════════════════════════════════════════════════

\subsection{Limited Time Span (31 Days)}

The dataset covers only January 2021. This limits:
\begin{itemize}
    \item \textbf{Seasonal analysis}: We cannot detect monthly or yearly
          cycles. Weekly cycles are the longest detectable pattern.
    \item \textbf{Model training}: 31 data points (daily) is small for
          complex models. SARIMA and simple regression are appropriate;
          deep learning (LSTM) is not, unless we aggregate at finer
          granularity (hourly: $31 \times 24 = 744$ points).
    \item \textbf{Forecasting horizon}: With 31 days total and an 80/20
          split, we forecast only $\sim$6 days. This is sufficient for
          demonstration but limits real-world applicability.
\end{itemize}

\subsection{Data Granularity}

The data has \textbf{minute-level resolution} (Date + Time). This enables
session analysis and fine-grained vibe assignment but also means aggregation
choices matter:
\begin{itemize}
    \item \textbf{Hourly aggregation}: 744 points, good for diversity and
          session metrics.
    \item \textbf{Daily aggregation}: 31 points, good for SARIMA and
          high-level trends.
    \item \textbf{Weekly aggregation}: $\sim$4 points, too few for ML but
          good for summary statistics.
\end{itemize}

The ideal approach is to compute metrics at \textbf{daily resolution} for
time series modelling, and at \textbf{hourly resolution} when more data
points are needed (e.g.\ for neural models or gap analysis).

\subsection{Missing Genre Information}

The dataset contains Artist, Track, and Album but \textbf{no genre or audio
features}. This means:
\begin{itemize}
    \item We cannot directly compute ``vibe'' in the musical sense (e.g.\
          energy, valence, danceability).
    \item However, genre can be \textbf{approximated} by using artist-level
          clusters from SVD on the user-artist interaction matrix (as
          explored in the \texttt{svd\_kmeans\_test/} directory).
    \item External APIs (Last.fm tags, Spotify audio features) could enrich
          the dataset if desired.
\end{itemize}

% ══════════════════════════════════════════════════════════════════════════════
\section{Conclusion}
% ══════════════════════════════════════════════════════════════════════════════

Raw play count is merely the surface of what \texttt{fm\_data.csv} can
reveal. By reframing the data through the lens of \textbf{vibes},
\textbf{diversity}, \textbf{discovery}, \textbf{sessions}, and
\textbf{loyalty}, we unlock a rich multi-dimensional portrait of listening
behaviour.

The ``vibes'' framework is the most significant conceptual contribution:
it transforms timestamps from raw numbers into meaningful behavioural
categories, enabling both human-interpretable summaries and powerful
machine learning features. When crossed with metrics like Shannon entropy
or discovery rate, vibes reveal \emph{why} a user listens differently at
different times --- not just \emph{that} they do.

Every metric explored in this document is computable from the existing data,
produces a meaningful time series plot, and has a well-defined machine
learning modelling strategy. The recommended priority order balances
novelty, visual impact, computational feasibility, and insight depth.

The next step is to implement these analyses, starting with the vibe
framework and diversity index, building toward a comprehensive dashboard
that integrates all perspectives into a unified view of music listening
behaviour.

% ══════════════════════════════════════════════════════════════════════════════
\end{document}
